% Wersja źródłowa: audyt III, jednostka 05
\setcounter{section}{4}
\section{Przestrzenie lokalnie wypukłe i seminormy}\label{sec:lokalnie-wypukle}

\subsection{Motywacja: od topologii do nierówności}

W poprzednich rozdziałach badaliśmy przestrzenie liniowo-topologiczne
bez wymagania, aby ich otoczenia były wypukłe. Okazało się, że nawet
zupełna i metryzowalna TVS może zachowywać się odmiennie od przestrzeni
unormowanej: w $L^p([0,1])$, $0<p<1$, otoczka wypukła ograniczonej
kuli nie jest ograniczona, a każdy ciągły funkcjonał liniowy jest zerowy.
Ten drugi fakt wynika z następującego argumentu. Niech
$L:L^p([0,1])\to\K$ będzie ciągłym funkcjonałem liniowym.
Dla $f\in L^p$ piszemy
$\|f\|_p=(\int_{[0,1]}|f|^p\,\mathrm d\lambda)^{1/p}$;
przy $p<1$ nie jest to norma. Z ciągłości $L$ w zerze istnieje $\delta>0$ takie, że
$\|g\|_p<\delta$ implikuje $|L(g)|<1$.
Dla $g\ne0$ podstawiamy $\delta g/(2\|g\|_p)$,
a z jednorodności dostajemy
$|L(g)|\le 2\|g\|_p/\delta$; nierówność ta
jest oczywista także dla $g=0$.
Przyjmujemy zatem $C=2/\delta$.
Jeśli $f\ne0$, to nieatomową miarę
$\nu_f(E)=\int_E|f|^p\,\mathrm d\lambda$ można dla
każdego $n\ge1$ podzielić na $n$ rozłącznych
zbiorów mierzalnych $E_1,\ldots,E_n$ o jednakowej
mierze $\nu_f(E_j)=\nu_f([0,1])/n$. Zatem
$f=\sum_{j=1}^n f\chi_{E_j}$ prawie wszędzie, a
\[
 |L(f)|\le\sum_{j=1}^n|L(f\chi_{E_j})|
 \le C n^{1-1/p}\|f\|_p\longrightarrow0.
\]
Ponieważ $1-1/p<0$, mamy $L(f)=0$ dla każdego $f$.
Wprowadzenie \emph{lokalnej wypukłości} pozwoli nam powiązać topologię
z oszacowaniami przez seminormy, a następnie z teorią dualności.

W całym rozdziale $\K\in\{\R,\C\}$. Jeżeli nie zaznaczono inaczej,
$X$ i $Y$ oznaczają przestrzenie liniowo-topologiczne nad $\K$;
nie zakładamy milcząco ani normowalności, ani metryzowalności.
Przez $\mathcal N_X(0)$ oznaczamy rodzinę otoczeń zera w $X$.
Zbiór $A\subseteq X$ nazywamy \emph{pochłaniającym}, jeśli
$\forall x\in X\ \exists t>0:\ x\in tA$.
W TVS każde otoczenie zera jest pochłaniające, ponieważ
$\lambda x\to0$ dla $\lambda\to0$.

\begin{definicja}[Wypukłość i zrównoważenie]
Niech $X$ będzie przestrzenią liniową nad $\K$.
Zbiór $A\subseteq X$ nazywamy \emph{wypukłym}, gdy
$(1-t)x+ty\in A$ dla wszystkich $x,y\in A$, $t\in[0,1]$.
Nazywamy go \emph{zrównoważonym}, gdy $\lambda A\subseteq A$
dla każdej $\lambda\in\K$ z $|\lambda|\le1$.
Zbiór jednocześnie wypukły i zrównoważony nazywamy
\emph{absolutnie wypukłym}.
\end{definicja}

\begin{stwierdzenie}[Wypukłość absolutna jako nierówność na sumę]
\label{prop:absolutna}
Niech $X$ będzie przestrzenią liniową nad $\K$ i
$A\subseteq X$ zbiorem niepustym. Jest on absolutnie wypukły
wtedy i tylko wtedy, gdy dla wszystkich $x,y\in A$ oraz
$\alpha,\beta\in\K$ spełniających $|\alpha|+|\beta|\le1$
zachodzi $\alpha x+\beta y\in A$.
\end{stwierdzenie}
\begin{proof}
Jeśli $A$ jest absolutnie wypukły, to
$\alpha x+\beta y=(|\alpha|+|\beta|)z$ przy
$z=\frac{|\alpha|}{|\alpha|+|\beta|}\frac{\alpha}{|\alpha|}x+
\frac{|\beta|}{|\alpha|+|\beta|}\frac{\beta}{|\beta|}y$
(wyrazy z zerowym współczynnikiem pomijamy).
Wypukłość i zrównoważenie dają $z\in A$ oraz
$\alpha x+\beta y\in A$. Gdy oba współczynniki są zerowe,
$0\in A$ wynika ze zrównoważenia. Odwrotnie, kładąc
$\beta=0$ otrzymujemy zrównoważenie, a kładąc
$\alpha=t$, $\beta=1-t$ dla $t\in[0,1]$, wypukłość.
\end{proof}

\begin{definicja}[Przestrzeń lokalnie wypukła]
\label{def:lcs}
Przestrzeń liniowo-topologiczną $X$ nazywamy
\emph{lokalnie wypukłą} (LCS), jeśli zero ma bazę otoczeń
złożoną ze zbiorów wypukłych. Przez \emph{lokalnie wypukłą
przestrzeń Hausdorffa} rozumiemy LCS, której topologia jest
Hausdorffa; rozróżnienie to jest istotne w twierdzeniach o normowalności.
\end{definicja}

\begin{stwierdzenie}[Absolutnie wypukła baza otoczeń]
\label{prop:absconv-base}
Niech $X$ będzie TVS. Przestrzeń $X$ jest lokalnie wypukła
wtedy i tylko wtedy, gdy zero ma bazę otoczeń absolutnie wypukłych.
Bazę można wybrać ze zbiorów otwartych.
\end{stwierdzenie}
\begin{proof}
Załóżmy lokalną wypukłość i ustalmy otoczenie $U$ zera.
Wybierzmy wypukłe otoczenie $V\subseteq U$, a następnie
otwarte zrównoważone otoczenie $O\subseteq V$.
Zbiór $\operatorname{conv}(O)$ jest zawarty w $V$,
bo $V$ jest wypukły; jest też zrównoważony,
ponieważ dla $|\lambda|\le1$ zachodzi
$\lambda\operatorname{conv}(O)=\operatorname{conv}(\lambda O)
\subseteq\operatorname{conv}(O)$.
Zbiór $\operatorname{conv}(O)$ jest absolutnie wypukły.
Ponadto jest otwarty: jeżeli $x\in\operatorname{conv}(O)$,
to $x+O\subseteq\operatorname{conv}(O)+O\subseteq
2\operatorname{conv}(O)$ nie wystarcza do żądanej tezy;
można natomiast zapisać $x=\sum_{j=1}^m t_jx_j$
z $t_j\ge0$, $\sum t_j=1$, i wybrać $j$ z $t_j>0$.
Ponieważ $O$ jest otwarty, dla dostatecznie małego
otoczenia $N$ zera mamy $x_j+t_j^{-1}N\subseteq O$,
stąd $x+N\subseteq\operatorname{conv}(O)$.
Druga implikacja wynika bezpośrednio z definicji.
\end{proof}

\begin{uwaga}[Dlaczego lokalna wypukłość jest istotna?]
W przestrzeni $L^p([0,1])$ dla $0<p<1$ nie ma bazy
wypukłych otoczeń zera; dowód podano w
przykładzie z \hyperref[sec:rodziny-metryzowalnosc]{rozdziału 4}.
Tak więc zupełność i metryzowalność nie wymuszają lokalnej wypukłości.
\end{uwaga}

\subsection{Seminormy i topologie zadane przez rodziny seminorm}

\begin{definicja}[Seminorma]
Niech $X$ będzie przestrzenią liniową nad $\K$.
\emph{Seminormą} na $X$ nazywamy funkcję
$p:X\to[0,\infty)$ spełniającą
$p(\lambda x)=|\lambda|p(x)$ dla $\lambda\in\K$, $x\in X$,
oraz $p(x+y)\le p(x)+p(y)$ dla $x,y\in X$.
Jeżeli dodatkowo $p(x)=0$ implikuje $x=0$, to $p$ jest normą.
\end{definicja}

\begin{przyklad}[Pierwsze seminormy]
Na $\K^{\N}$ funkcje $p_n(x)=|x_n|$ są seminormami, ale
nie normami. Na przestrzeni $C^\infty([0,1])$ seminormami są
$p_m(f)=\max_{0\le j\le m}\sup_{t\in[0,1]}|f^{(j)}(t)|$,
$m\in\N\cup\{0\}$. Na $C(\Omega)$ dla otwartego
$\Omega\subseteq\R^d$ seminormą jest
$p_K(f)=\sup_{x\in K}|f(x)|$ dla dowolnego zwartego $K\subseteq\Omega$.
\end{przyklad}

\begin{twierdzenie}[Topologia generowana przez seminormy]
\label{thm:seminorm-topology}
Niech $X$ będzie przestrzenią liniową nad $\K$, a
$\mathcal P$ niepustą rodziną seminorm na $X$. Istnieje
najmniejsza topologia $\tau_{\mathcal P}$ \emph{zgodna ze
strukturą przestrzeni liniowo-topologicznej}, w której
wszystkie seminormy $p\in\mathcal P$ są ciągłe.
Topologia ta czyni $X$ przestrzenią lokalnie wypukłą,
a bazą otoczeń zera są zbiory
\[
 U(p_1,\ldots,p_m;\varepsilon)
 =\{x\in X:\max_{1\le j\le m}p_j(x)<\varepsilon\},
 \quad m\ge1,\quad p_j\in\mathcal P,\quad\varepsilon>0.
\]
Jest ona Hausdorffa wtedy i tylko wtedy, gdy
$\bigcap_{p\in\mathcal P}\ker p=\{0\}$.
Ponadto sieć $(x_d)$ zbiega do $x$ w $\tau_{\mathcal P}$
dokładnie wtedy, gdy $p(x_d-x)\to0$ dla każdego $p\in\mathcal P$.
\end{twierdzenie}
\begin{proof}
Zbiory z tezy są absolutnie wypukłe, pochłaniające i
zamknięte na skończone przecięcia (po ewentualnym
zmniejszeniu $\varepsilon$). Dla zbioru z promieniem
$\varepsilon$ suma dwóch analogicznych zbiorów z promieniem
$\varepsilon/2$ mieści się w wyjściowym zbiorze.
Aby uzasadnić ciągłość mnożenia przez skalar w punkcie
$(\lambda_0,x_0)$, dla każdego z ustalonych $p_j$
używamy oszacowania
$p_j(\lambda x-\lambda_0x_0)\le
|\lambda|p_j(x-x_0)+|\lambda-\lambda_0|p_j(x_0)$.
Wystarczy ograniczyć $|\lambda|$ w małym otoczeniu
$\lambda_0$, a następnie wybrać odpowiednio małe
$|\lambda-\lambda_0|$ i $p_j(x-x_0)$ dla wszystkich
skończenie wielu indeksów $j$.
W ten sposób bazy przesuniętych zbiorów definiują
topologię TVS.
Każde $p\in\mathcal P$ jest ciągłe w zerze, a z nierówności
$|p(x)-p(y)|\le p(x-y)$ wynika jej ciągłość wszędzie.
Każda inna topologia TVS, w której seminormy są ciągłe,
zawiera wszystkie powyższe otoczenia, zatem jest silniejsza.

Jeżeli $x\ne0$ i $p(x)>0$ dla pewnego $p$, otoczenie
$\{y:p(y)<p(x)/2\}$ oddziela $0$ od $x$.
Jeśli natomiast $p(x)=0$ dla wszystkich $p$, to $x$
należy do każdego otoczenia zera. W konsekwencji
$\tau_{\mathcal P}$ jest Hausdorffa dokładnie przy
podanym warunku. Opis zbieżności wynika z bazy otoczeń.
\end{proof}

\subsection{Funkcjonał Minkowskiego}

\begin{definicja}[Funkcjonał Minkowskiego]
\label{def:gauge}
Niech $X$ będzie przestrzenią liniową, a $A\subseteq X$
niepustym zbiorem pochłaniającym. Funkcjonałem Minkowskiego
(\emph{funkcjonałem cechującym}) zbioru $A$ nazywamy
\[
 p_A(x)=\inf\{t>0:x\in tA\},\qquad x\in X.
\]
Pochłanianie zapewnia, że infimum jest skończone.
\end{definicja}

\begin{twierdzenie}[Od zbioru absolutnie wypukłego do seminormy]
\label{thm:minkowski}
Niech $X$ będzie TVS nad $\K$, a $A\subseteq X$
pochłaniającym zbiorem absolutnie wypukłym.
Wówczas $p_A$ jest seminormą i
$\{x:p_A(x)<1\}\subseteq A\subseteq\{x:p_A(x)\le1\}$.
Jeśli dodatkowo $A$ jest \emph{otwarty}, to
$A=\{x:p_A(x)<1\}$ i $p_A$ jest ciągła.
\end{twierdzenie}
\begin{proof}
Pochłanianie oraz $0\in A$ dają $p_A(x)<\infty$ i
$p_A(0)=0$. Dla $\lambda\ne0$ mamy
$x\in tA$ dokładnie wtedy, gdy $\lambda x\in|\lambda|tA$:
przy zespolonym $\lambda$ korzystamy ze zrównoważenia
$A$ i z równości $\frac{\lambda}{|\lambda|}A=A$.
Otrzymujemy $p_A(\lambda x)=|\lambda|p_A(x)$;
przypadek $\lambda=0$ jest osobny i natychmiastowy.

Niech $s>p_A(x)$ oraz $t>p_A(y)$. Z definicji infimum
i z własności $rA\subseteq sA$ dla $0<r\le s$
wynika $x\in sA$ i $y\in tA$. Wypukłość daje
$sA+tA\subseteq(s+t)A$, zatem
$p_A(x+y)\le s+t$. Przechodząc z $s$ i $t$ do
odpowiednich infimów otrzymujemy podaddytywność.
Jeżeli $p_A(x)<1$, to $x\in tA$ dla pewnego $t<1$,
więc $x\in A$ dzięki zrównoważeniu.
Jeżeli $x\in A$, to $p_A(x)\le1$.

Niech teraz $A$ będzie otwarty i niech $x\in A$.
Ciągłość $\lambda\mapsto\lambda x$ w $\lambda=1$
daje $\lambda>1$ dostatecznie bliskie $1$ z
$\lambda x\in A$; stąd $p_A(x)\le1/\lambda<1$.
Zatem $A=\{p_A<1\}$.
Ponieważ $\{p_A<\varepsilon\}=\varepsilon A$
dla każdego $\varepsilon>0$, seminorma jest ciągła
w zerze, a nierówność
$|p_A(x)-p_A(y)|\le p_A(x-y)$ zapewnia ciągłość wszędzie.
\end{proof}

\begin{center}
\begin{tikzpicture}[scale=.85,>=Latex]
 \draw[fill=blue!5,draw=blue!50!black,thick] (0,0) ellipse (2.3 and 1.35);
 \draw[->,gray] (-2.8,0)--(4.35,0) node[right] {$x_1$};
 \draw[->,gray] (0,-1.65)--(0,1.9) node[above] {$x_2$};
 \fill (0,0) circle (1.8pt) node[below left] {$0$};
 \draw[->,thick] (0,0)--(3.45,1.12) node[above right] {$x$};
 \fill[blue!70!black] (2.012,.653) circle (2.3pt);
 \node[blue!60!black,above] at (1.50,.84) {$x/p_A(x)$};
 \node[blue!60!black] at (-1.45,.7) {$A$};
\end{tikzpicture}

\small Rysunek: dla absolutnie wypukłego otwartego otoczenia $A$
dla punktu $x$ z $p_A(x)>0$ wektor $x/p_A(x)$ leży
na brzegu $A$; wartość $p_A(x)$ określa współczynnik skalowania.
\end{center}

\begin{przyklad}[Seminorma nie musi być normą]
Dla $X=\R^2$ i $A=\{(x,y):|x|<1\}$ zachodzi
$p_A(x,y)=|x|$. Zbiór $A$ jest otwarty, absolutnie
wypukły i pochłaniający, lecz nie jest ograniczony,
a $p_A(0,1)=0$. Takie zdegenerowanie nie przeczy
Twierdzeniu~\ref{thm:minkowski}.
\end{przyklad}

\begin{przyklad}[Pochłanianie nie gwarantuje ciągłości]
Niech $X=c_{00}$ z normą $\|x\|_\infty$ oraz
$A=\{x\in c_{00}:\sup_{n\ge1}n|x_n|\le1\}$.
Każdy $x\in c_{00}$ ma skończony nośnik, więc $A$
jest pochłaniający i absolutnie wypukły. Jego funkcjonał
to $p_A(x)=\sup_{n\ge1}n|x_n|$. Nie jest on ciągły
w topologii normy supremum: $x^{(n)}=e_n/n\to0$,
ale $p_A(x^{(n)})=1$. W szczególności $A$ nie jest
otoczeniem zera w wyjściowej topologii.
\end{przyklad}

\subsection{Charakteryzacja lokalnej wypukłości przez seminormy}

\begin{twierdzenie}[Charakteryzacja lokalnej wypukłości]
\label{thm:lcs-seminorms}
Dla dowolnej TVS $X$ nad $\K$ następujące warunki są
równoważne: (i) $X$ jest lokalnie wypukła;
(ii) topologia $X$ jest generowana przez rodzinę
ciągłych seminorm. W przypadku Hausdorffa rodzinę
można wybrać rozdzielającą punkty.
\end{twierdzenie}
\begin{proof}
Implikację (ii)$\Rightarrow$(i) daje
Twierdzenie~\ref{thm:seminorm-topology}.
Załóżmy (i) i weźmy rodzinę $\mathcal U$ wszystkich
otwartych absolutnie wypukłych otoczeń zera.
Dla każdego $U\in\mathcal U$ określmy funkcjonał
Minkowskiego $p_U$ z poprzedniej podsekcji. Dowód
Twierdzenia~\ref{thm:minkowski} pokazuje, że $p_U$
jest ciągłą seminormą i $U=\{p_U<1\}$. Zatem
rodzina $(p_U)_{U\in\mathcal U}$ generuje wyjściową
topologię. Jeśli $X$ jest Hausdorffa, dla $x\ne0$
możemy wybrać $U\in\mathcal U$ niezawierające $x$,
a wówczas $p_U(x)\ge1$. Żaden niezerowy wektor
nie należy do wszystkich jąder.
\end{proof}


\subsection{Ciągłość odwzorowań liniowych i ograniczoność}

\begin{twierdzenie}[Kryterium ciągłości przez seminormy]
\label{thm:continuous-sem}
Niech $X,Y$ będą przestrzeniami lokalnie wypukłymi,
$\mathcal P$ i $\mathcal Q$ rodzinami seminorm generującymi
odpowiednio ich topologie, a $T:X\to Y$ odwzorowaniem liniowym.
Następujące warunki są równoważne:
\begin{enumerate}[label=\textup{(\roman*)}]
\item $T$ jest ciągłe;
\item dla każdej $q\in\mathcal Q$ istnieją $m\ge1$,
$p_1,\ldots,p_m\in\mathcal P$ i $C>0$ takie, że
$q(Tx)\le C\max_{1\le j\le m}p_j(x)$ dla wszystkich $x\in X$.
\end{enumerate}
\end{twierdzenie}
\begin{proof}
Załóżmy ciągłość $T$. Dla $q\in\mathcal Q$
znajdźmy $p_1,\ldots,p_m$ i $\delta>0$ takie, że
$\max_jp_j(x)<\delta$ implikuje $q(Tx)<1$.
Oznaczmy $r(x)=\max_jp_j(x)$. Jeżeli $r(x)=0$,
to $r(tx)=0$ dla każdego $t>0$, a więc
$tq(Tx)<1$ dla każdego $t>0$; stąd $q(Tx)=0$.
Jeżeli $r(x)>0$, to dla
$t=\delta/(2r(x))$ mamy $r(tx)=\delta/2$,
a zatem $q(Tx)<2r(x)/\delta$. Uzyskujemy (ii)
z $C=2/\delta$. Odwrotnie, każda z nierówności
z (ii) zapewnia ciągłość $q\circ T$ w zerze;
opis bazy otoczeń z Twierdzenia~\ref{thm:seminorm-topology}
w $Y$ daje ciągłość $T$.
\end{proof}

\begin{twierdzenie}[Ograniczoność przez seminormy]
\label{thm:bounded-sem}
Niech $X$ będzie przestrzenią lokalnie wypukłą,
a $\mathcal P$ rodziną seminorm generującą topologię $X$.
Zbiór $B\subseteq X$ jest ograniczony w sensie von Neumanna
wtedy i tylko wtedy, gdy
$\sup_{x\in B}p(x)<\infty$ dla każdego $p\in\mathcal P$
(przyjmujemy $\sup\varnothing=0$).
\end{twierdzenie}
\begin{proof}
Jeśli $B$ jest ograniczony, to dla danego $p$ i
$V=\{x:p(x)<1\}$ istnieje $t>0$ z $B\subseteq tV$,
więc $\sup_Bp\le t$.
Odwrotnie, niech $U$ będzie otoczeniem zera.
Wybierzmy $p_1,\ldots,p_m\in\mathcal P$ i
$\varepsilon>0$ z
$\{x:\max_jp_j(x)<\varepsilon\}\subseteq U$.
Niech $M=\max_j\sup_{x\in B}p_j(x)<\infty$.
Dla $t>(M/\varepsilon)$ mamy
$B\subseteq t\{x:\max_jp_j(x)<\varepsilon\}\subseteq tU$.
\end{proof}

\begin{wniosek}[Otoczki ograniczone]
\label{cor:absconv-bounded}
Jeżeli $X$ jest przestrzenią lokalnie wypukłą,
a $B\subseteq X$ jest ograniczony, to jego otoczka wypukła
i jego otoczka absolutnie wypukła są ograniczone.
\end{wniosek}
\begin{proof}
Dla każdej ciągłej seminormy $p$ i dla kombinacji
$z=\sum_{j=1}^m\lambda_jx_j$ z $x_j\in B$ oraz
$\sum_j|\lambda_j|\le1$ zachodzi
$p(z)\le\sup_{x\in B}p(x)$. Takie kombinacje
wypełniają otoczkę absolutnie wypukłą $B$.
Stosujemy Twierdzenie~\ref{thm:bounded-sem}.
\end{proof}

\subsection{Twierdzenie Kolmogorowa o normowalności}

\begin{twierdzenie}[Kryterium normowalności]
\label{thm:kolmogorov-normability}
Niech $X$ będzie lokalnie wypukłą przestrzenią
liniowo-topologiczną Hausdorffa nad $\K$.
Następujące warunki są równoważne:
\begin{enumerate}[label=\textup{(\roman*)}]
\item topologia $X$ pochodzi od pewnej normy na $X$;
\item $X$ ma ograniczone otoczenie zera;
\item $X$ ma ograniczone, otwarte i absolutnie wypukłe otoczenie zera.
\end{enumerate}
\end{twierdzenie}
\begin{proof}
Jeżeli $X$ jest unormowana, jej otwarta kula jednostkowa
jest otoczeniem ograniczonym. To daje (i)$\Rightarrow$(ii).
Dzięki lokalnej wypukłości wewnątrz ograniczonego
otoczenia możemy wybrać otwarte absolutnie wypukłe
otoczenie $U$, które jest ograniczone jako podzbiór
zbioru ograniczonego. Mamy (ii)$\Rightarrow$(iii).

Załóżmy (iii). Funkcjonał Minkowskiego $p_U$ jest
ciągłą seminormą i $U=\{p_U<1\}$.
Pokażemy, że jest normą. Jeśli $p_U(x)=0$,
to dla każdego $s>0$ zachodzi $x\in sU$.
Niech $V$ będzie dowolnym otoczeniem zera w $X$.
Ograniczoność $U$ daje $t>0$ z $U\subseteq tV$,
zatem $x\in t^{-1}U\subseteq V$.
Ponieważ $X$ jest Hausdorffa, przecięcie wszystkich
otoczeń zera jest równe $\{0\}$, czyli $x=0$.

Topologia normy $p_U$ nie jest silniejsza od wyjściowej,
ponieważ $p_U$ jest ciągła. Odwrotnie, dla dowolnego
otoczenia $V$ ograniczoność $U$ daje
$t^{-1}U\subseteq V$, a zbiór ten jest kulą
$\{p_U<1/t\}$. Zatem obie topologie są równe.
\end{proof}

\begin{uwaga}[Gdzie konieczna jest własność Hausdorffa?]
Jeżeli lokalna wypukłość jest spełniona, ale przestrzeń
nie jest Hausdorffa, funkcjonał Minkowskiego ograniczonego
otoczenia może być tylko seminormą z niezerowym jądrem.
Topologia generowana przez seminormę nieoddzielającą punktów
nie może pochodzić od normy.
\end{uwaga}

\subsection{Metryzowalność lokalnie wypukła i przestrzenie Frécheta}

\begin{twierdzenie}[Przeliczalna rodzina seminorm]
\label{thm:countable-seminorms}
Niech $X$ będzie lokalnie wypukłą TVS Hausdorffa.
Jest ona metryzowalna wtedy i tylko wtedy, gdy jej topologię
generuje przeliczalna rodzina seminorm $(p_n)_{n\ge1}$.
W takim przypadku zgodną, niezmienniczą względem przesunięć
metryką jest
\[
 d(x,y)=\sum_{n=1}^{\infty}2^{-n}\min\{1,p_n(x-y)\}.
\]
\end{twierdzenie}
\begin{proof}
Jeżeli istnieje przeliczalna rodzina $(p_n)$,
to każda funkcja $\min\{1,p_n(x-y)\}$ spełnia
nierówność trójkąta, a szereg jest jednostajnie zbieżny
(bo jego wyrazy nie przekraczają $2^{-n}$).
Ponieważ seminormy rozdzielają punkty, $d$ jest metryką.
Jeśli $x_k\to x$ w topologii seminorm, to dla ustalonego
$N$ pierwszych $N$ składników sumy dążą do zera,
a ogon jest jednostajnie mniejszy od $2^{-N}$;
stąd $d(x_k,x)\to0$. W drugą stronę z
$d(x_k,x)\ge2^{-n}\min\{1,p_n(x_k-x)\}$
wynika zbieżność każdej seminormy. Ten sam argument
zamiast ciągów stosuje się do otoczeń i dowodzi
równości topologii.

Jeżeli $X$ jest metryzowalna, posiada przeliczalną bazę
otoczeń zera (kryterium metryzowalności TVS z \hyperref[sec:rodziny-metryzowalnosc]{rozdziału 4}).
Do każdego jej elementu wybierzmy, korzystając z
Twierdzenia~\ref{thm:lcs-seminorms}, skończoną rodzinę
ciągłych seminorm i wspólny promień tak, aby odpowiadający
zbiór był zawarty w danym otoczeniu. Suma tych skończonych
rodzin jest przeliczalna i generuje topologię $X$.
\end{proof}

\begin{definicja}[Przestrzeń Frécheta]
\label{def:frechet-space}
\emph{Przestrzenią Frécheta} nazywamy przestrzeń
liniowo-topologiczną, która jest lokalnie wypukła,
metryzowalna, zupełna i Hausdorffa. Nie wymagamy, by
metryka generująca topologię była zadana normą.
\end{definicja}

\begin{przyklad}[Produkt, który nie jest normowalny]
\label{ex:product-lcs}
Niech $X=\K^{\N}$ z topologią produktową.
Generują ją seminormy $p_n(x)=|x_n|$.
Topologia jest Hausdorffa, metryzowalna, a przestrzeń
jest zupełna jako produkt zupełnych przestrzeni $\K$.
Jest zatem przestrzenią Frécheta. Nie jest natomiast
normowalna: żadne otoczenie zera nie jest ograniczone,
co pokazaliśmy w \hyperref[sec:rodziny-metryzowalnosc]{rozdziale 4}.
\end{przyklad}

\begin{przyklad}[Gładkie funkcje i nieskończenie wiele pochodnych]
\label{ex:smooth-frechet-preview}
Niech $X=C^\infty([0,1])$ z topologią generowaną przez
$p_m(f)=\max_{0\le j\le m}\sup_{t\in[0,1]}|f^{(j)}(t)|$,
$m=0,1,\ldots$. Seminormy te rozdzielają punkty,
a topologia jest metryzowalna i lokalnie wypukła.
Zupełność wykażemy w następnym stwierdzeniu. Już teraz
można dowieść \emph{nienormowalności}. Gdyby $U$ było ograniczonym
otoczeniem zera, zawierałoby
$V=\{f:p_m(f)<\delta\}$ dla pewnych $m$ i $\delta>0$;
zbiór $V$ byłby więc ograniczony. Tymczasem dla
$N\ge2\pi$ funkcje
$f_N(t)=\frac{\delta}{2N^m}\sin(Nt)$
spełniają $p_m(f_N)\le\delta/2$, natomiast
$p_{m+1}(f_N)\ge\delta N/2\to\infty$.
To przeczy Twierdzeniu~\ref{thm:bounded-sem}.
Przestrzeń $C^\infty([0,1])$ jest więc lokalnie wypukła
oraz metryzowalna, ale nie normowalna.
\end{przyklad}


\subsection{Sieci Cauchy'ego i zupełność w języku seminorm}

\begin{twierdzenie}[Kryterium Cauchy'ego w przestrzeniach lokalnie wypukłych]
\label{thm:lcs-cauchy}
Niech $X$ będzie przestrzenią lokalnie wypukłą, a
$\mathcal P$ rodziną seminorm generującą topologię $X$.
Sieć $(x_d)_{d\in D}$ w $X$ jest siecią Cauchy'ego
w kanonicznej strukturze jednostajnej TVS dokładnie wtedy,
gdy dla każdej $p\in\mathcal P$ i każdego $\varepsilon>0$
istnieje $d_0\in D$ takie, że $p(x_d-x_e)<\varepsilon$
dla wszystkich $d,e\succeq d_0$.
W szczególności $X$ jest zupełna dokładnie wtedy, gdy każda
taka sieć jest zbieżna w $X$.
\end{twierdzenie}
\begin{proof}
Jeżeli sieć jest Cauchy'ego, zastosujmy definicję z
\hyperref[sec:zupelnosc]{rozdziału 3} do otoczenia $\{x:p(x)<\varepsilon\}$.
Otrzymujemy żądany warunek dla każdej seminormy.
Na odwrót, niech $U$ będzie dowolnym otoczeniem zera.
Znajdźmy seminormy $p_1,\ldots,p_m$ i $\varepsilon>0$
takie, że $\{x:\max_jp_j(x)<\varepsilon\}\subseteq U$.
Dla każdego $j$ mamy indeks $d_j$, po którym
$p_j(x_d-x_e)<\varepsilon$; skierowanie $D$ daje wspólny
indeks $d_0\succeq d_j$, $j=1,\ldots,m$.
Wtedy $x_d-x_e\in U$ dla $d,e\succeq d_0$.
Ostatnie zdanie jest dokładnie definicją zupełności TVS.
\end{proof}

\begin{twierdzenie}[Zupełność przestrzeni funkcji gładkich]
\label{thm:smooth-complete}
Niech $X=C^\infty([0,1])$ nad $\K$ z topologią generowaną
przez seminormy $p_m(f)=\max_{0\le j\le m}
\sup_{t\in[0,1]}|f^{(j)}(t)|$, $m=0,1,\ldots$.
Wtedy $X$ jest przestrzenią Frécheta, ale nie jest normowalna.
\end{twierdzenie}
\begin{proof}
Lokalna wypukłość, warunek Hausdorffa, metryzowalność
i brak normowalności zostały wykazane powyżej.
Dla zupełności, na mocy twierdzenia o metryzowalnych
TVS z \hyperref[sec:zupelnosc]{rozdziału 3}, wystarczy wykazać zupełność ciągową.

Niech $(f_n)$ będzie ciągiem Cauchy'ego względem wszystkich
$p_m$. Dla każdego $j\ge0$ ciąg funkcji ciągłych
$(f_n^{(j)})$ jest jednostajnie Cauchy'ego na $[0,1]$.
Zupełność $C([0,1])$ w normie supremum daje
funkcję ciągłą $g_j$ taką, że $f_n^{(j)}\to g_j$
jednostajnie. Korzystając z podstawowego twierdzenia
rachunku całkowego, dla $t\in[0,1]$ mamy
\[
 f_n^{(j)}(t)-f_n^{(j)}(0)
 =\int_{[0,t]}f_n^{(j+1)}(s)\,\mathrm d\lambda(s).
\]
Jednostajna zbieżność całkowanych funkcji i wartości
w zerze pozwala przejść do granicy:
$g_j(t)-g_j(0)=\int_{[0,t]}g_{j+1}(s)\,\mathrm d\lambda(s)$.
Ponieważ $g_{j+1}$ jest ciągła, wynika stąd
$g_j'=g_{j+1}$. Indukcyjnie $g_0$ jest gładka i
$g_0^{(j)}=g_j$ dla wszystkich $j\ge0$.
Ponadto $p_m(f_n-g_0)\to0$ dla każdego $m$.
Oznacza to, że ciąg Cauchy'ego ma granicę w $X$,
a więc $X$ jest zupełna.
\end{proof}

\begin{uwaga}[Rzut na skończenie wiele pochodnych]
Z punktu widzenia konstrukcji przestrzeni funkcyjnych
$C^\infty([0,1])$ można uważać za zgodną rodzinę
funkcji należących do przestrzeni Banacha $C^m([0,1])$,
$m\ge0$, z normami $p_m$. Każdy operator
$f\mapsto f^{(j)}(t_0)$, $t_0\in[0,1]$, jest ciągłym
funkcjonałem liniowym na tej przestrzeni, gdyż
$|f^{(j)}(t_0)|\le p_j(f)$. Tę strukturę będziemy
wykorzystywać przy definiowaniu dystrybucji.
\end{uwaga}

\subsection{Podprzestrzenie, produkty i ilorazy}

\begin{stwierdzenie}[Trwałość lokalnej wypukłości]
\label{prop:lcs-constructions}
Niech $(X_i)_{i\in I}$ będzie rodziną przestrzeni lokalnie
wypukłych nad $\K$. Wówczas:
\begin{enumerate}[label=\textup{(\roman*)}]
\item każda podprzestrzeń liniowa $M\subseteq X_i$ z topologią
podprzestrzeni jest lokalnie wypukła;
\item produkt $\prod_{i\in I}X_i$ z topologią produktową
jest lokalnie wypukły;
\item dla dowolnej podprzestrzeni liniowej $M\subseteq X_i$
iloraz $X_i/M$ z topologią ilorazową jest lokalnie wypukły.
Iloraz jest Hausdorffa wtedy i tylko wtedy, gdy $M$ jest domknięta.
\end{enumerate}
\end{stwierdzenie}
\begin{proof}
W (i) przecinamy wypukłe otoczenia zera w $X_i$
z $M$. W (ii) podstawowe otoczenia zera są produktami,
w których jedynie skończenie wiele współrzędnych
jest ograniczonych do wypukłych otoczeń; takie produkty
są wypukłe. W (iii) naturalne rzutowanie
$q:X_i\to X_i/M$ jest otwarte, a obrazy wypukłych
otoczeń zera są wypukłymi otoczeniami zera w ilorazie.
Kryterium Hausdorffa udowodniono w
\hyperref[sec:konstrukcje]{rozdziale 2}.
\end{proof}

\begin{uwaga}[Dalszy kierunek]
Związek topologii z seminormami jest punktem wyjścia
teorii dualności. W następnym rozdziale sformułujemy i
udowodnimy twierdzenie Hahna--Banacha, które pozwoli
przejść od seminorm i rozdzielania punktów do istnienia
nietrywialnych ciągłych funkcjonałów liniowych.
Później wprowadzimy topologie słabe, przestrzenie Frécheta,
przestrzenie funkcji testowych $\mathcal D(\Omega)$ oraz ich
przestrzeń sprzężoną $\mathcal D'(\Omega)$.
\end{uwaga}
